\section{ALNS clásico: selección por ruleta adaptativa}
\label{sec:clasico}

\cod{ALNS\_Classic} implementa el mecanismo adaptativo de Ropke y
Pisinger~\cite{ropke2006}. Mantiene dos ruletas \emph{independientes}, una para los $6$
operadores de destrucción y otra para los $5$ de reparación. Cada ruleta guarda, para cada
operador $i$, un peso $w_i$ (inicialmente $1$), un puntaje acumulado $S_i$ y un contador de
usos $u_i$ dentro del segmento en curso.

\subsection{Selección}

En cada iteración se sortea un operador de cada ruleta con probabilidad proporcional a su peso:
\begin{equation}
P(i)=\frac{w_i}{\sum_{l} w_l}.
\label{eq:ruleta}
\end{equation}
Como los dos sorteos son independientes, la probabilidad de un par es el producto
$P^{-}(j)\,P^{+}(k)$: el método no puede representar que un operador de destrucción funcione
mejor con un operador de reparación concreto.

\subsection{Asignación de crédito y actualización}

Tras evaluar la candidata, ambos operadores del par reciben el mismo puntaje $\sigma$, que
depende únicamente de la \emph{categoría} del resultado y no de su magnitud:
\begin{equation}
\sigma=
\begin{cases}
33 & \text{si la candidata es un nuevo mejor global,}\\
13 & \text{si mejora a la solución actual,}\\
9  & \text{si no mejora pero es aceptada,}\\
0  & \text{si es rechazada.}
\end{cases}
\label{eq:sigma}
\end{equation}
Los tres valores positivos son los propuestos en~\cite{ropke2006}, asignados aquí en orden
monótono con la calidad del resultado. Las categorías se determinan comparando valores de
$f$, por lo que valen igual en las dos fases: en la fase~1, colocar un cliente del banco es
una mejora (y normalmente un nuevo mejor de la etapa), y una candidata que deja más clientes
sin asignar que la solución actual es, en la práctica, un rechazo. Los pesos no cambian en cada iteración sino al cierre
de cada segmento de $\ell=100$ iteraciones, con una media móvil exponencial sobre el puntaje
medio del segmento:
\begin{equation}
w_i\gets\max\Bigl(w_{\min},\ \lambda\,w_i+(1-\lambda)\,\frac{S_i}{u_i}\Bigr)\quad\text{si }u_i>0,
\qquad \lambda=0.9,\quad w_{\min}=0.01.
\label{eq:pesos}
\end{equation}
Un operador que no se usó en el segmento conserva su peso. El Algoritmo~\ref{alg:clasico}
reúne ambas operaciones.

\begin{algorithm}[H]
\caption{Selección y aprendizaje del ALNS clásico (\cod{ALNS\_Classic})}
\label{alg:clasico}
\SetKwProg{Proc}{procedimiento}{}{}
\Proc{\textnormal{\fn{Seleccionar}($t$)}}{
  $j\gets$ sorteo en la ruleta de destrucción según~\eqref{eq:ruleta}\;
  $k\gets$ sorteo en la ruleta de reparación según~\eqref{eq:ruleta}\;
  \Devolver $(j,k)$\;
}
\Proc{\textnormal{\fn{Aprender}($t$, $o$)}}{
  $\sigma\gets$ puntaje de $o$ según~\eqref{eq:sigma}\;
  $S^{-}_j\gets S^{-}_j+\sigma$;\quad $u^{-}_j\gets u^{-}_j+1$\;
  $S^{+}_k\gets S^{+}_k+\sigma$;\quad $u^{+}_k\gets u^{+}_k+1$\;
  \Si(\Com*[f]{cierre de segmento}){$t \bmod 100=0$}{
    \ParaCada{ruleta y cada operador $i$ de la ruleta}{
      \lSi{$u_i>0$}{$w_i\gets\max\bigl(w_{\min},\ \lambda w_i+(1-\lambda)S_i/u_i\bigr)$}
      $S_i\gets 0$;\quad $u_i\gets 0$\;
    }
  }
}
\end{algorithm}

\subsection{Propiedades relevantes}

\begin{itemize}
\item \textbf{Normalización por uso.} El peso se actualiza con el puntaje \emph{medio}
$S_i/u_i$, por lo que un operador elegido con frecuencia no se refuerza solo por haber sido
usado más veces.
\item \textbf{Inercia.} Con $\lambda=0.9$ cada segmento aporta el 10\,\% del peso nuevo. El
puntaje medio está acotado en $[0,33]$, de modo que los pesos convergen hacia ese rango de
forma gradual: partiendo de $w_i=1$, hacen falta varios segmentos (varios cientos de
iteraciones) para que la distribución se aleje de la uniforme. Los segmentos se cuentan
sobre el índice global de iteración, sin distinguir etapas ni fases: en las $25\,000$
iteraciones de la fase~2 hay $250$ actualizaciones de pesos, a las que se suman hasta otras
$250$ en la fase~1.
\item \textbf{Continuidad entre fases.} El método no reacciona al cambio de fase
(\cod{onDistancePhase} conserva su implementación vacía): los pesos aprendidos mientras se
eliminaban rutas son el punto de partida de la fase de distancia, y se readaptan al nuevo
objetivo solo a través de la media móvil~\eqref{eq:pesos}.
\item \textbf{Exploración permanente.} La selección es proporcional, no voraz, y $w_{\min}$
impide que un peso se anule. Todos los operadores conservan una probabilidad no nula durante
toda la búsqueda, y la distribución se adapta a la fase en curso (por ejemplo, los operadores
que aceptan empeoramientos puntúan más cuando la temperatura es alta).
\item \textbf{Crédito ordinal.} Una mejora mínima de distancia y la colocación de un cliente
del banco reciben el mismo puntaje si pertenecen a la misma categoría; el método mide con qué frecuencia
un operador tiene éxito, no cuánto mejora.
\item \textbf{Sin contexto.} La distribución no depende del estado de la búsqueda más que a
través de los pesos acumulados.
\end{itemize}
